\section{The Metaphysics of Law}

\textbf{Valentin Tomberg} wrote about the nature of law in his two doctoral dissertations, \textit{Foundations of International Law as Humanity's Law} and \textit{Degeneration and Regeneration of Jurisprudence}. What follows is a summary of some aspects of those studies.

A complete understanding of law must be based on three insights:

\begin{enumerate}


\item The world is an ordered whole

\item The world order is of a moral nature

\item Mankind is called to arrange its own life in harmony with the moral world order
\end{enumerate}


This follows from the understanding that the world is ordered by the Logos, it is a creation of the Will, and the Will must be directed by the Spirit.

In terms of the schema described in The Degrees of Knowledge\footnote{\url{https://gornahoor.net/?p=411}}, the study of law may be done on three levels, as shown in Table \ref{tab:ThreeLevelsLaw}.

\begin{table}[h]\small\centering
\begin{tabular}{cccp{.25\textwidth}}\toprule

Doxa & Positive Law & Concept of Law & The study of the actual laws of various nations and times.\\ 
Dianoia & Natural Law & Idea of Law & The understanding that the purpose of law is justice \\
Episteme & Divine Law & Ideal of Law & Direct intuition of the divine basis of law \\\bottomrule
\end{tabular}
\caption{Three levels of Law}
\label{tab:ThreeLevelsLaw}
\end{table}

Tomberg summarizes:

\begin{quotex}
Thus one arrives at the \textbf{ideal of law} by introspection, at the \textbf{idea of law} by reasonable evaluation of experiences, and at the \textbf{concept of law} by knowing the facts of positive law.
\end{quotex}


A complete understanding of the law requires all three levels. A legal positivist, for example, will understand a ``right'' as whatever the government legislates rather than acknowledging its divine origin. Or someone may argue on the rational level in favour of economic or social justice, which are really based on the idea of equality, rather than justice.

Every positive law must be both \textbf{just} and \textbf{possible}. A just law entails two conditions:

\begin{itemize}


\item It must not contradict a higher obligation

\item It must be necessary or useful for the common welfare
\end{itemize}


This means a law cannot contradict the natural or divine law. Furthermore, it rules out frivolous, arbitrary, and pointless laws, if they don't serve the common welfare. Examples from the past include bans on playing card games or the period of Prohibition in the USA.

A law must also be physically and morally possible. There may be otherwise just laws that are above the common moral level of the general public, so that their enforcement would disrupt the civil order. So despite the principles that ``error has no rights'' and ``doing wrong is not a right'', certain behaviours may need to be tolerated. In other words, not every moral wrong needs to be treated as a legal matter. It is important, however, that mere toleration doesn't devolve into being considered a right.

Since part of Tomberg's thesis deals with the proper education of lawyers, there is the tendency to think that simply the study of these two essays should be sufficient to transform the curricula of law schools and legislatures themselves. The problem is that it will only be understood at the level of doxa and dianoia, and not at the level of episteme. That will relegate it to the role of just one more legal theory to be debated, rather than as the ordering principle of legal theory itself. In a Traditional society, men incapable of grasping ideas directly through episteme would not even be allowed to teach at or attend a law school. They would find other roles in society appropriate to their intellectual capacities.

Tomberg's theses may be obtained from the Sophia Foundation\footnote{\url{http://sophiafoundation.org/activities/hermeticism/study_materials}}.


\flrightit{Posted on 2010-05-03 by Cologero}

